\section{Introduction}

The design of a pharmaceutical manufacturing facility must ensure not only technical and economic feasibility but also the highest standards of safety, product quality, and environmental protection. This report evaluates the safety and environmental considerations associated with the proposed industrial production of semaglutide, a glucagon-like peptide-1 (GLP-1) receptor agonist widely used for the treatment of type-2 diabetes and obesity.

Semaglutide is a complex peptide manufactured through multi-step chemical synthesis and purification processes involving organic solvents and reactive reagents. As a result, the facility must operate under strict pharmaceutical quality standards while managing chemical hazards, occupational exposures, and waste streams generated during production. Failures in safety management or environmental protection within pharmaceutical manufacturing have historically led to severe consequences for patients, workers, and surrounding communities. Therefore, understanding lessons from past industrial incidents and integrating them into plant design is essential for preventing similar failures.

This report applies the principles of Best Available Techniques (BAT) \cite{EU_IED_2010_75} and inherently safer design (ISD) \cite{Yang_ISD_2018} to minimise hazards at the source rather than relying solely on downstream controls. BAT involves adopting the most effective and advanced methods that are technically and economically viable to reduce environmental impact, emissions, and waste generation. Inherently safer design principles (minimisation, substitution, moderation, and simplification) are incorporated where possible to reduce the quantity of hazardous materials handled, replace hazardous substances with safer alternatives, and simplify process operations to limit the potential for accidents.

Moreover, applicable safety and environmental legislation governing pharmaceutical production are reviewed. Based on this context, structured risk assessment methods such as qualitative risk assessment and Hazard and Operability Study (HAZOP) are used to identify potential process hazards and inform appropriate safeguards, mechanical integrity considerations, and site layout decisions. The environmental impact of the plant is also assessed through evaluation of raw material selection, waste disposal routes, as well as utilities and greenhouse gas emissions. Overall, this section informs the development of safe operating practices and effective waste management strategies to ensure the facility operates safely, sustainably, and in compliance with modern pharmaceutical manufacturing standards.
